\section{Introduction}
\IEEEPARstart{H}{olding} a bus can improve service spacing while delaying the passengers already on board. Under a fixed fleet, the same seconds can also postpone the vehicle's next departure. A terminal timetable revision therefore depends not only on demand and the current terminal queue, but also on how much vehicle time downstream holding will consume. This study examines that connection through a shared cumulative holding budget.

The service--holding tradeoff is established in the bus-control literature. Daganzo's headway-based control seeks reliable high-frequency service with reduced schedule slack~\cite{daganzo2009headway}. Daganzo and Pilachowski extend cooperation to front and rear spacing, using speed adjustment to counter large gaps~\cite{daganzo2011cooperation}. Berrebi \emph{et al.} compare predictive and non-predictive holding methods in terms of headway stability and holding time~\cite{berrebi2018comparing}. These findings motivate a direct question: does a turnover-aware allocation rule improve the operating tradeoff relative to simple holding restrictions, or does it mainly choose a different point on that tradeoff?

An executable comparison requires consistent resource semantics. A timetable predictor can assume a returning vehicle will be available while later holding delays that return. A cumulative allowance can also have different effects from a per-stop cap even when both lead to similar average holding. We address these issues by using the same published releases, irreversible commands, and remaining-credit ledger inside each forecast and in physical execution. This makes a budget an inspectable resource constraint rather than an additional reward term.

The architecture also permits a controlled examination of temporal information and learning. Causal harmonic summaries inform a bounded timetable law; the fast layer evaluates legal holding actions on current automatic passenger count (APC) and automatic vehicle location (AVL) observations. A learned supervisor can enable or withhold timetable revisions, but the same interface can be driven by an always-joint manual supervisor. Matching their budget settings separates a learning increment from a change in holding exposure. Matching a harmonic and non-harmonic representation addresses a different question again.

We present three contributions. First, we specify the time-budget interface and its accounting invariants under a fixed fleet and a fixed trip roster. Second, independent paired-demand experiments establish how retention changes passenger exposure, circulation, and actual regularity, while separating those effects from supervisor choice. Third, strong simple-limit comparisons and a common-reference admission redesign expose the present boundary of pressure-driven allocation. The resulting evidence identifies both an operationally useful configuration effect and the unresolved challenge of improving it without transferring cost to service regularity.

\section{Related Work}
\subsection{Holding, Prediction, and Vehicle Circulation}
Headway-based regulation uses spacing information to control unstable service~\cite{daganzo2009headway}. Cooperative control can additionally exploit the state of the following bus~\cite{daganzo2011cooperation}. The native continuation used here is a discrete bilateral holding rule; it is not the continuous speed controller of Daganzo and Pilachowski. Its availability provides a defined fallback when the event model lacks support.

Prediction-based holding introduces a further distinction between the action preferred under a forecast and its realized service effect. The predictive/non-predictive comparison of Berrebi \emph{et al.}~\cite{berrebi2018comparing} is particularly relevant to assessing headway stability alongside the intervention time required. Our comparison adds a fixed physical vehicle pool, public departure obligations, and a cumulative allowance that can affect both return forecasts and executable commands. The tests retain simple credit and stop-cap controls rather than interpreting improvement over an old learner as sufficient evidence.

Asynchronous multi-agent reinforcement learning addresses the non-simultaneous decisions made when different buses reach stops~\cite{wang2021asynchronous}. Our fast decisions are also event-triggered, but are selected by finite model-based enumeration. The optional learned component instead acts on the slow publication clock. This distinction matters for interpreting which component is responsible for an observed operating change.

\subsection{Temporal Representation and Action Persistence}
Multiple control frequencies can be formalized through action persistence, including policies whose action components operate on different clocks~\cite{lee2020multiple}. Persistence can itself be adapted or optimized in batch reinforcement learning~\cite{metelli2020persistence}. A persistent holding allowance is a resource commitment across stop events; it is not repeated execution of one fixed holding command.

Similarly, a harmonic representation of observations is not equivalent to frequency-specific action authority. Here, harmonic summaries enter the slow state and timetable law. The fast layer uses APC clearance evidence, AVL, plans, and remaining credit. The representation comparison therefore tests a defined change to this implementation, not a complete separation of low- and high-frequency controllers. The learned comparator uses a finite-action maximum-entropy actor--critic~\cite{haarnoja2018sac,christodoulou2019discrete}; its optimizer is retained as an experimental component rather than a claimed methodological advance.
